% vim:fileencoding=utf-8:foldmethod=marker:foldmarker=(((,)))

% \documentclass[8pt]{extarticle}  % COMENTAME AL TERMINAR
% \usepackage{geometry}            % COMENTAME AL TERMINAR
% \geometry{paperwidth=8.5cm,      % COMENTAME AL TERMINAR
					% paperheight=10.8cm,    % COMENTAME AL TERMINAR
					% margin=0.6cm}          % COMENTAME AL TERMINAR
% \pagestyle{empty}                % COMENTAME AL TERMINAR
% \usepackage{amsmath}             % COMENTAME AL TERMINAR
% \usepackage{amssymb}             % COMENTAME AL TERMINAR
% \usepackage{mathrsfs}            % COMENTAME AL TERMINAR
% \pagestyle{empty}                % COMENTAME AL TERMINAR
% \setlength{\columnsep}{1pt}      % COMENTAME AL TERMINAR
% \begin{document}                 % COMENTAME AL TERMINAR

% --- PÁGINA 3 ---

% === HEADER === (((
\begin{center}
    {\large \textbf{ $\S$ Jardín de Estructuras}} {$~$\small\FourClowerSolid }
\end{center}
\vspace{0.14cm}
\begin{center}
    {--- \large\textit{ Espacios Topológicos }---}
\end{center}
\vspace{0.15cm}
% )))

% === THE CONTENT === (((
{\fontsize{7}{5} Recordemos que si $X$ es un conjunto... 
Entre algunos ejemplos tenemos:}
\begin{center}
\footnotesize 
\setlength{\tabcolsep}{4pt}
\begin{tabular}{@{} c c m{2.55cm} @{}}
\hline 
\textbf{Espacio} & \textbf{Notación} & \textbf{Descripción}\\
\hline 
Discreto & $(X, \mathscr{P}(X))$ & \raggedright\arraybackslash\scriptsize Metrizable y $0$-dimensional. \\[4pt]
Indiscreto & $(X, \{\varnothing, X\})$ & \raggedright\arraybackslash\scriptsize Pseudometrizable y compacto. \\
\hline
\end{tabular}
\end{center}
% )))

% \end{document}                   % COMENTAME AL TERMINAR
